% Figure: a ball bouncing with coefficient of restitution e < 1. Each
% bounce returns to e^2 times the previous height, so the peaks form a
% geometric sequence and the whole motion ends in a finite time.
\begin{figure}[htbp]
\centering
\begin{tikzpicture}[scale=1.0]
% ground
\draw[engblack, thick] (0,0) -- (11.2,0);
% first drop from h0 = 3, then peaks at e^2 h with e^2 = 0.64
% peak heights: 3, 1.92, 1.229, 0.786, 0.503
% We draw parabolic arcs between contact points spaced by shrinking widths.
% contact x-positions chosen to shrink like the time between bounces (factor e)
\def\hzero{3.0}
% arc 1: fall from (0,hzero) to (1.6,0)
\draw[engblue, very thick] (0,\hzero) .. controls (0.8,1.3) .. (1.6,0);
% arc 2: up to peak 1.92 then down; width scales by e=0.8 -> 2.56 wide
\draw[engblue, very thick] (1.6,0) .. controls (2.88,2.56) .. (4.16,0);
% arc 3: peak 1.229, width 2.048
\draw[engblue, very thick] (4.16,0) .. controls (5.184,1.64) .. (6.208,0);
% arc 4: peak 0.786, width 1.638
\draw[engblue, very thick] (6.208,0) .. controls (7.027,1.05) .. (7.846,0);
% arc 5: peak 0.503, width 1.311
\draw[engblue, very thick] (7.846,0) .. controls (8.501,0.67) .. (9.157,0);
% remaining bounces collapse toward the accumulation point
\draw[engblue, very thick] (9.157,0) .. controls (9.68,0.43) .. (10.2,0);
\draw[engblue, very thick] (10.2,0) .. controls (10.53,0.27) .. (10.86,0);
% accumulation point marker
\fill[engred] (11.0,0) circle (2pt);
\node[font=\scriptsize, engred, anchor=south west] at (11.02,0.02)
  {rest at finite $t_\infty$};
% height markers
\draw[enggray, dashed] (0,\hzero) -- (-0.15,\hzero);
\node[font=\scriptsize, anchor=east] at (-0.15,\hzero) {$h_0$};
\draw[enggray, dashed] (2.88,1.92) -- (-0.15,1.92);
\node[font=\scriptsize, anchor=east] at (-0.15,1.92) {$e^{2}h_0$};
\draw[enggray, dashed] (5.184,1.229) -- (-0.15,1.229);
\node[font=\scriptsize, anchor=east] at (-0.15,1.229) {$e^{4}h_0$};
% axis label
\node[font=\scriptsize, anchor=north] at (5.6,-0.15) {time / horizontal offset (drawn shrinking by $e$ each bounce)};
\end{tikzpicture}
\caption[Bounce heights of a ball with restitution below one]{A ball
dropped from height $h_0$ onto a floor of restitution $e<1$ returns to
$e^{2}h_0$ after the first bounce, $e^{4}h_0$ after the second, and
$e^{2n}h_0$ after the $n$th, a geometric sequence of peaks. The time
between successive bounces shrinks by the factor $e$, so the arcs
crowd together and the infinite sequence of bounces completes in a
finite total time $t_\infty$. The ball comes to rest after infinitely
many ever-shorter bounces, the same accumulation a mathematician meets
in a convergent geometric series and an engineer hears as the run of
clicks that ends in silence.}
\label{fig:vol03ch04:bounce-series}
\end{figure}
